% !TEX root = ../PolycopiePremiere2627.tex
%%%%%%%%%%%%%%%%%%%%%%%%%%%%%%%%%%%%%%%%%%%%%%%%%%%%%%%%%%%%%%%%%%%%%%%%%%%%%%%%
% Contenu initial repris de Louis Paternault --- http://ababsurdo.fr
%
% Publié sous licence Creative Commons Attribution-ShareAlike 4.0 International (CC BY-SA 4.0)
% http://creativecommons.org/licenses/by-sa/4.0/deed.fr
%%%%%%%%%%%%%%%%%%%%%%%%%%%%%%%%%%%%%%%%%%%%%%%%%%%%%%%%%%%%%%%%%%%%%%%%%%%%%%%%

\chapter{Probabilités conditionnelles}

\section{Fréquences marginales, fréquences conditionnelles}

\Dfn{Dans un tableau croisé d'effectifs :
\begin{itemize}
    \item la valeur à l'intersection d'une ligne et d'une colonne est le nombre d'individus présentant simultanément les deux caractères ;
    \item on appelle \textcolor{J1}{effectif marginal} l'effectif d'une sous-population qui ne dépend que d'un seul caractère. Ces effectifs sont représentés dans la ligne « Total » ou la colonne « Total » ;
    \item la \textcolor{J1}{fréquence marginale} d'une valeur est le quotient :
    \[\dfrac{\text{Effectif marginal}}{\text{Effectif total}}\]
    \item la \textcolor{J1}{fréquence conditionnelle} de la valeur $x_i$ sachant $y_j$ est le quotient :
    \[\dfrac{\text{Effectif présentant les valeurs $x_i$ et $y_j$}}{\text{Effectif marginal de la valeur $y_j$}}\]
\end{itemize}}

\Exemple{Un nouveau test a été développé pour détecter une maladie, et on souhaite évaluer son efficacité. On prend un échantillon représentatif de \numprint{1000} personnes dans la population, dont on sait par ailleurs si elles sont malades ou non, et on leur fait passer le test. On obtient les résultats suivants (où « diagnostic positif » signifie que le test a détecté la maladie, peut-être à tort).
\begin{center}
    \begin{tabular}{|M{3.4cm}|M{1.8cm}|M{1.8cm}|M{1.8cm}|}
        \hline
        & \textbf{Malade} & \textbf{Saine} & \textbf{Total} \tabularnewline
        \hline
        Diagnostic positif & 19 & 79 & \tabularnewline
        \hline
        Diagnostic négatif & & 901 & \tabularnewline
        \hline
        Total & & & \tabularnewline
        \hline
    \end{tabular}
\end{center}

\begin{enumerate}
    \item Complétez le tableau.
    \item Combien de personnes malades ont obtenu un test négatif ?
    \item Calculez la fréquence marginale de personnes malades dans la population.
    \item Calculez la fréquence de diagnostics corrects. Le test est-il fiable ?
    \item Calculez la fréquence conditionnelle de personnes malades sachant que le diagnostic est positif.
\end{enumerate}}

\section{Probabilités conditionnelles}

\Dfn[Rappels]{Étant donnés deux évènements $A$ et $B$ d'un univers $\Omega$ :
\begin{itemize}
    \item l'intersection de $A$ et $B$, notée $A\cap B$, est \textcolor{J1}{l'ensemble des éléments qui appartiennent à la fois à $A$ et à $B$} ;
    \item l'union de $A$ et $B$, notée $A\cup B$, est \textcolor{J1}{l'ensemble des éléments qui appartiennent à $A$, ou à $B$, ou aux deux à la fois} ;
    \item le contraire de $A$, noté $\overline{A}$, est \textcolor{J1}{l'ensemble des issues de l'univers qui n'appartiennent pas à $A$}.
\end{itemize}}

\Exemple{On considère les élèves de cette classe, et :
\begin{itemize}
    \item $A$ : « l'ensemble des garçons » ;
    \item $B$ : « l'ensemble des personnes portant des lunettes ».
\end{itemize}
Alors :
\begin{itemize}
    \item $A\cup B$ est : \textcolor{J1}{l'ensemble des garçons à lunettes}.
    \item $\overline{B}$ est : \textcolor{J1}{l'ensemble des personnes sans lunettes}.
    \item $A\cap\overline{B}$ est : \textcolor{J1}{l'ensemble des garçons sans lunettes}.
\end{itemize}}

\Dfn[Rappel]{On choisit de manière équiprobable un individu au hasard dans une population, et on considère un évènement $A$. Alors :
\[
    P(A)=\textcolor{J1}{\dfrac{\text{Nombre de cas favorables}}{\text{Nombre de cas possibles}}}
\]}

\Exemple{Pour vérifier la qualité d'une chaîne de montage d'aspirateurs, une équipe vérifie leur fonctionnement : sur 67 aspirateurs étudiés, 3 ne fonctionnent pas.

On choisit un aspirateur au hasard à la sortie de cette chaîne de montage. Quelle est la probabilité qu'il ne fonctionne pas ?}

\Dfn{Soient $A$ et $B$ deux évènements d'une expérience aléatoire, tels que $P(B)\neq0$.

On appelle \textcolor{J1}{probabilité conditionnelle de $A$ sachant $B$}, ou « probabilité que $A$ soit réalisé sachant que $B$ est réalisé », le nombre :
\[P_B\left( A \right)=\dfrac{\mathrm{card}(A\cap B)}{\mathrm{card}(B)}=\dfrac{P\left( A\cap B \right)}{P\left( B \right)}\]}

\Exemple{On lance un dé équilibré à six faces, et on observe le résultat obtenu. On note $P$ : « le résultat est pair », et $Q$ : « le résultat est supérieur ou égal à quatre ».
\begin{enumerate}
    \item Calculer la probabilité que le résultat soit pair.
    \item Calculer la probabilité que le résultat soit pair, sachant qu'il est supérieur ou égal à 4.
\end{enumerate}}

\Dfn{Deux évènements $A$ et $B$ (avec $P(B)\neq0$) sont dits \textcolor{J1}{indépendants} si $P_B(A)=P(A)$.}

\Prp{Deux évènements $A$ et $B$ sont \textcolor{J1}{indépendants} si et seulement si $P(A\cap B)=P(A)\times P(B)$.}

\Exemple[Inspiré du sujet 0 du baccalauréat d'enseignement spécifique, sujet 2, 2025]{
% Dactylographié par l'APMEP. Merci à elles et eux.
Dans un lycée comptant $\numprint{2000}$ élèves, on donne la répartition des effectifs suivant le sexe et le choix de la LV1.

\begin{center}
    \begin{tabular}{|M{2.6cm}|M{2cm}|M{2cm}|}
        \hline
        & \textbf{Fille} & \textbf{Garçon} \tabularnewline
        \hline
        Anglais & 712 & 728 \tabularnewline
        \hline
        Autre LV1 & 288 & 272 \tabularnewline
        \hline
    \end{tabular}
\end{center}

\begin{enumerate}
    \item Un élève affirme « Dans ce lycée, il y a autant de filles que de garçons ». A-t-il raison ? Justifier.
\end{enumerate}

On choisit au hasard, de manière équiprobable, un élève dans ce lycée.

On considère les évènements suivants :

$F$ : « l'élève est une fille » ;

$A$ : « l'élève a choisi Anglais pour LV1 ».

Dans les questions qui suivent, on donnera les résultats sous forme d'une fraction qu'il n'est pas demandé de simplifier.

\begin{enumerate}[resume]
    \item Déterminer la probabilité de l'évènement $A \cap F$.
    \item Déterminer la probabilité de l'évènement $A$ sachant que l'évènement $F$ est réalisé.
    \item Les évènements $A$ et $F$ sont-ils indépendants ? Justifier.
    \item On sait que l'élève choisi est un garçon. On considère l'affirmation suivante :

    « La probabilité qu'il ait choisi Anglais pour LV1 est plus de trois fois plus grande que la probabilité qu'il n'ait pas choisi Anglais pour LV1 ».

    Cette affirmation est-elle vraie ? Justifier.
\end{enumerate}}

\pagebreak

\section{Répétition d'expériences aléatoires}

\Prp{On peut modéliser la succession d'expériences aléatoires par un arbre de probabilités (ou arbre pondéré) respectant les règles suivantes.
\begin{itemize}
    \item La somme des probabilités des branches issues d'un nœud est égale à 1.
    \item La probabilité d'une issue d'un chemin est égale au produit des probabilités rencontrées sur le chemin.
    \item La probabilité d'un évènement est égale à la somme des probabilités des issues qui le composent.
\end{itemize}}

\Exemple[Inspiré du sujet 0 du baccalauréat d'enseignement spécifique, sujet 3, 2025]{
% Dactylographié par Denis Vergès, de l'APMEP. Merci à lui.
Un village propose aux participants de la fête du sport deux épreuves : une randonnée et un cross. Il n'est pas possible de s'inscrire aux deux épreuves à la fois.
On dispose des informations suivantes :
\begin{itemize}
    \item $90\,\%$ des participants ont choisi la randonnée, parmi eux, $5\,\%$ sont licenciés dans un club ;
    \item $10\,\%$ des participants ont choisi le cross, parmi eux, $40\,\%$ sont licenciés dans un club.
\end{itemize}

Un journaliste interroge un participant au hasard, et on considère les évènements suivants :

$R$ : « Le participant a choisi la randonnée » ;

$L$ : « Le participant est licencié dans un club ».

\begin{enumerate}
    \item Par simple lecture de l'énoncé, indiquer :
    \begin{enumerate}
        \item La probabilité que le participant interrogé soit licencié dans un club sachant qu'il a choisi la randonnée.
        \item La probabilité que le participant interrogé soit licencié dans un club sachant qu'il a choisi le cross.
    \end{enumerate}
\end{enumerate}

\begin{enumerate}[resume]
    \item Représenter la situation par un arbre de probabilité.
    \item
    \begin{enumerate}
        \item Déterminer la probabilité que le participant interrogé ait choisi le cross et soit licencié dans un club.
        \item Vérifier que la probabilité que le participant interrogé soit licencié dans un club est égale à $\dfrac{850}{10000}$, soit $8,5\,\%$.
    \end{enumerate}
    \item Le journaliste interroge un participant licencié dans un club. Déterminer la probabilité que ce participant ait choisi le cross.
\end{enumerate}}
